\section{Introduction}\label{sec:introduction}

Computing the roots of a polynomial by solving a companion eigenvalue problem
changes the perturbations to which the answer is exposed. A polynomial root
responds to errors in its coefficients; the unstructured eigenvalue condition
permits errors in every entry of the companion matrix. We ask how the conditioning of these two problems
compares for Gaussian input. Our main result gives a sharp probabilistic
description of companion eigenvalue conditioning and a relative comparison
with the original root problem.

Average-condition theorems for matrix polynomials use their coefficient
matrices as input. Applying such a theorem after linearization also requires
its probability law and perturbation norm to match those of the induced
companion problem. Section~\ref{sec:intro-previous} makes this distinction
explicit for the results of Armentano--Beltr\'an.

Consider
\[
 p(x)=\sum_{k=0}^n a_kx^k,
 \qquad a_0,\ldots,a_n\ \text{independent standard complex Gaussians}.
\]
Here each coefficient has independent real and imaginary parts of variance
$1/2$. Let $C_p$ be the $n\times n$ monic companion matrix, with ones on its
subdiagonal and last column $-(a_0,\ldots,a_{n-1})^T/a_n$. First sample $p$,
then choose one of its $n$ roots $Z$ uniformly; write $\Pstar,\Estar$ for this
joint law. Almost surely the roots are simple and nonzero.

For a root $z$, let $r,w$ be right and left eigenvectors of $C_p$ at $z$:
$C_pr=zr$ and $w^*C_p=zw^*$. The quantity
\begin{equation}\label{eq:intro-overlap}
 \Lambda(p,z)=\frac{\|w\|_2\|r\|_2}{|w^*r|}
\end{equation}
is the absolute eigenvalue condition number for arbitrary matrix perturbations
in spectral norm. It is often called the eigenvector overlap. We compare
relative condition numbers using Euclidean norm on polynomial coefficients,
spectral norm on matrices, and relative error in the same scalar root. Their
precise definitions are given in Section~\ref{sec:setup}.

\Needspace{17\baselineskip}
\begin{theorem}[Main result: conditioning of Gaussian companions]\label{thm:intro-main}
There is an absolute constant $C$ such that the following hold.
\begin{enumerate}[(i)]
 \item For every degree $n\ge1$ and every $t>0$,
 \begin{equation}\label{eq:intro-tail}
   \Pstar\{\Lambda(p,Z)>t\}\le\min\{1,Ct^{-2}\}.
 \end{equation}
 Moreover,
 \begin{equation}\label{eq:intro-second}
   \Estar\Lambda(p,Z)^2=2\log n+O(1)
   \qquad(n\to\infty).
 \end{equation}
 Consequently the moments $\Estar\Lambda^q$ are bounded uniformly in degree
 precisely for $0<q<2$, and $\Estar\log\Lambda=O(1)$.
 \item The relative eigenvalue and coefficient condition numbers satisfy
 \begin{equation}\label{eq:intro-relative}
   \Estar\log
   \frac{\kappa_{\mathrm{eig,rel}}(C_p,Z)}
        {\kappa_{\mathrm{coeff,rel}}(p,Z)}
   =\log n+O(1)
   \qquad(n\to\infty).
 \end{equation}
\end{enumerate}
The constants in the $O(1)$ terms are independent of $n$.
\end{theorem}

The first assertion controls the expected fraction of roots with a large
eigenvalue condition number, uniformly in the degree. Its exponent is sharp:
any larger exponent with a degree-independent constant would bound the second
moment, contradicting \eqref{eq:intro-second}. The second assertion says that
the geometric mean of the relative companion-to-coefficient condition ratio
is comparable to $n$. These statements fit together because the coefficient
root condition has its own degree scale. They concern a uniformly selected
root and two specified local problems, rather than a bound on all errors in
a complete rootfinding algorithm.

For clarity, the two relative conditions in the theorem are
\begin{equation}\label{eq:intro-relative-definitions}
 \kappa_{\mathrm{coeff,rel}}(p,z)
 =\frac{\|a\|_2\bigl(\sum_{k=0}^n|z|^{2k}\bigr)^{1/2}}
        {|z|\,|p'(z)|},\qquad
 \kappa_{\mathrm{eig,rel}}(C_p,z)
 =\Lambda(p,z)\frac{\|C_p\|_2}{|z|},
\end{equation}
where $a=(a_0,\ldots,a_n)^T$ and every polynomial coefficient is perturbable.
Theorem~\ref{thm:comp-main} proves part (i), and
Corollary~\ref{cor:comp-relative} proves part (ii).

\subsection{Proof strategy and the role of root geometry}

The proof separates local root sensitivity from the additional conditioning
of interpolation in the companion basis. To analyze local sensitivity, we
first work with systems of $d$ independent complex Gaussian polynomials
sharing a covariance kernel. That kernel defines a natural output metric.
Marked Kac--Rice shows that, after normalization by this metric, the Jacobian
at a selected root is independent of the root location and has a
determinant-weighted Gaussian law. The resulting intrinsic condition number
$\kappa_H$, defined in Section~\ref{sec:setup-kernel-metric}, satisfies
\[
 \Pstar\{\kappa_H>t\}\le
 \min\left\{1,\frac{d(d+1)}{4t^4}\right\},
\]
with the constant asymptotically exact for fixed $d$. The precise assumptions
and statements are Theorems~\ref{thm:local-universal} and~\ref{thm:local-tail}.
The induced Ginibre and Wishart distributions used here are established
random-matrix results \cite{fischmann2012induced,kumar2017smallest}.

For the Gaussian box ensemble, exact radial laws then translate this geometric
normalization back to Euclidean root displacement. They give
$\Estar\log\kappa_{\mathrm{coeff}}=-\log n+O_d(1)$ and retain the
contribution of roots outside the unit polydisk. Finally, in one variable,
writing $p(x)=(x-z)q(x)$ turns the coefficient norm into the Gram matrix of
multiplication by $x-z$. This matrix is tridiagonal. Its inverse trace, together
with the Gaussian weight at a root, gives the exact conditional overlap moment
and the bounds in Theorem~\ref{thm:intro-main}. The same argument also shows
that, for each fixed $n\ge2$, overlap moments are finite precisely below order
four; the threshold two concerns uniformity as the degree grows.

The paper follows this proof structure. Section~\ref{sec:setup} fixes the
input--output maps, metrics, and sampling law.
Section~\ref{sec:local-law} establishes the local Gaussian law;
Section~\ref{sec:box} derives the box root and coefficient scales; and
Section~\ref{sec:companion} proves the main theorem.
Section~\ref{sec:discussion} discusses its interpretation and the multivariate
questions that remain.

\subsection{Relation to previous work}\label{sec:intro-previous}

Local conditioning of polynomial systems has a substantial geometric and
probabilistic theory. Shub--Smale developed the projective framework, its
relation to ill-posedness, and distributional estimates in the
\emph{Complexity of B\'ezout's theorem} series
\cite{shub1993geometric,shub1993volumes,shub1993packing}.
B\"urgisser--Cucker \cite{burgisser2013condition} give a systematic account.
Proposition~\ref{prop:setup-condition} is a standard derivative-norm
calculation in explicitly specified input and output metrics; it is included
to fix conventions, not as a new condition-number theorem.

Beltr\'an--Pardo \cite{beltran2009smale,beltran2011fast} used average
conditioning and integration over system--root pairs to obtain average
polynomial-time homotopy algorithms. Their work also gives exact moments of
the normalized projective condition number below order four
\cite[Sections~7--8]{beltran2011fast}.
B\"urgisser--Cucker \cite{burgisser2011smale} established smoothed estimates,
and Armentano--Beltr\'an--B\"urgisser--Cucker--Shub
\cite[Theorem~2]{armentano2016condition} obtained an exact centered
inverse-Jacobian Frobenius second moment and a smoothed bound.
Section~\ref{sec:local-classical} spells out the overlap with our local law.
Lairez \cite{lairez2017deterministic} subsequently obtained a deterministic
algorithm with polynomial average complexity.

For sparse systems, Malajovich--Rojas \cite{malajovich2004high} construct
metrics from the coefficient spaces, relate conditioning to the discriminant,
and prove fourth-power tail bounds. Their normalization and worst-root
probability differ from our absolute, selected-root formulation. Toric
condition metrics and homotopy complexity are developed further in
\cite{malajovich2019toric,malajovich2023renormalization}.
Our local section presents the common-kernel Gaussian calculation in a form
suited to comparing coefficient and algebraic-representation perturbations;
the fourth-power exponent and the homogeneous moment identities already
belong to this established theory.

Toh--Trefethen \cite{toh1994pseudozeros} explain the different perturbation
geometries of polynomial roots and companion eigenvalues. Noferini--Townsend
and Graf--Townsend construct multivariate examples in which intermediate
algebraic problems amplify conditioning exponentially with the number of
variables \cite{noferini2016numerical,graf2026numerical}. Our representation
theorem concerns one variable and growing degree; the local framework extends
to systems, but the present paper does not resolve those growing-dimension
examples.

B\"urgisser--Cucker--Rocha Cardozo \cite{burgisser2015characteristic} study
the conditioning of characteristic polynomials induced by Gaussian matrices.
Their results demonstrate the importance of the induced coefficient law in
an algebraic reduction. Here the construction runs from a Gaussian polynomial
to its companion matrix, giving a different probability distribution and
perturbation problem.

\Needspace{4\baselineskip}
Armentano--Beltr\'an \cite[Theorem~1.1]{armentano2019polynomial} study the
polynomial eigenvalue problem for
$Q(z)=\sum_{j=0}^k B_jz^j$, where $B_j\in\C^{m\times m}$.
Their perturbable input is the tuple $(B_0,\ldots,B_k)$, with norm
$(\sum_j\|B_j\|_F^2)^{1/2}$, and their random model makes all entries of these
coefficient matrices independent complex Gaussians. Thus their main theorem
concerns matrix-polynomial eigenvalue conditioning; Remark~1.4 discusses
companion algorithms in this setting as well as scalar polynomial rootfinding.
The scalar specialization \cite[Section~3.1]{armentano2019polynomial} gives an
exact identity for our coefficient ensemble, using the standard Fubini--Study
output metric. In our notation,
\[
 \mu_{\mathrm{FS}}(p,z)
 =\frac{\|a\|_2\bigl(\sum_{k=0}^n|z|^{2k}\bigr)^{1/2}}
        {(1+|z|^2)|p'(z)|},\qquad \Estar\mu_{\mathrm{FS}}^2=1.
\]
This is relative to the input norm and uses a projective output metric,
explained in Section~\ref{sec:setup-kernel-metric}. It differs from our kernel
metric. The companion overlap concerns a further change of input: arbitrary
perturbations of $C_p$. Viewing $zI-C_p$ as a degree-one matrix polynomial does
not make the Gaussian coefficient-matrix theorem applicable: its leading
matrix is fixed, and $C_p$ has prescribed zeros and ones and ratios of Gaussian
coefficients. Our theorem analyzes this induced matrix ensemble. The distinction
is therefore between input spaces, metrics, and probability laws, rather than
between polynomial solving and eigenvalue problems as such.

Barbarino--Noferini study limiting empirical eigenvalue measures of random
matrix polynomials, both with a fixed identity leading coefficient
\cite{barbarino2023monic} and with random leading coefficients
\cite{barbarino2022limit}. The latter paper includes our scalar Gaussian
ensemble and its limiting uniform root distribution on the unit circle
\cite[Theorem~1.2]{barbarino2022limit}. These results describe root locations;
weak convergence of empirical measures does not determine companion overlap
moments. Their monic paper also provides bounds for the smallest singular
value of a Gaussian matrix with arbitrary deterministic mean
\cite[Theorem~A.5]{barbarino2023monic}. These are relevant tools for smoothed
analysis, although an application to root conditioning must account for
sampling at a root and for the actual perturbation structure. Both papers
treat polynomials in one variable; growing matrix size is distinct from a
growing number of variables in a polynomial system.
